\documentclass[11pt]{article}

\usepackage[margin=1in]{geometry}
\usepackage{amsmath,amssymb,amsthm}
\usepackage{booktabs}
\usepackage{enumitem}
\usepackage[expansion=false]{microtype}
\usepackage{xcolor}
\usepackage[colorlinks=true,linkcolor=blue!55!black,citecolor=blue!55!black,
            urlcolor=blue!55!black]{hyperref}

\newtheorem{principle}{Design Principle}

\newcommand{\dist}{\operatorname{dist}}
\newcommand{\diam}{\operatorname{diam}}
\newcommand{\Sstar}{\mathcal{S}^\star}
\newcommand{\R}{\mathbb{R}}
\newcommand{\Z}{\mathbb{Z}}

\title{Geometric Heuristics for Basis Selection in the GSVD\\
Method of Particular Solutions}
\author{}
\date{}

\begin{document}

\maketitle

\begin{abstract}
\noindent
We give a closed-form heuristic system for choosing the type, count, and placement
of particular solutions in the GSVD formulation of the Method of Particular Solutions
(MPS) for the planar Dirichlet Laplace eigenproblem. The basis consists of
corner-adapted fractional-order Fourier--Bessel functions
$J_{k\alpha}(\sqrt{\lambda}\,r)\sin(k\alpha\theta)$ together with real-valued mono-
and multipole fundamental solutions $Y_k(\sqrt{\lambda}\,r)\{\sin,\cos\}(k\theta)$
placed at exterior source points. All prescriptions depend only on measurable
geometric data --- area, perimeter, diameter, corner angles, and inter-corner
distances --- and on a single accuracy parameter $\Lambda = \ln(1/\hat\epsilon)$
derived from the target eigenvalue precision via the Moler--Payne bound.
\end{abstract}

\tableofcontents

\bigskip

The system is built around two error mechanisms that behave completely differently
--- \emph{oscillatory resolution}, which scales like $\kappa=\sqrt{\lambda}$, and
\emph{analytic-continuation decay}, which scales like $\log(1/\epsilon)$ --- plus a
\emph{partition-of-responsibility} principle that decides which basis family owns
which piece of $\partial\Omega$. Everything below reduces to closed-form counts in
terms of $A=|\Omega|$, $L=|\partial\Omega|$, $D=\diam\Omega$, the corner angles
$\omega_c$, and inter-corner distances.

\section{Notation, and converting \texorpdfstring{$\epsilon$}{epsilon} into a target tension}
\label{sec:notation}

Let $\kappa=\sqrt{\lambda}$ and $\kappa_{\max}=\sqrt{\lambda_{\max}}$. Corner $c$ has
vertex $v_c$, interior angle $\omega_c$, and corner exponent
\[
  \alpha_c=\pi/\omega_c,
  \qquad
  \text{basis } J_{j\alpha_c}(\kappa r_c)\sin(j\alpha_c\theta_c),
  \quad j=1,\dots,M_c .
\]

The tension $t=\|u\|_{L^2(\partial\Omega)}/\|u\|_{L^2(\Omega)}$ has units of
length$^{-1/2}$, so form the dimensionless quantity
\[
  \hat t \;=\; \sqrt{A/L}\;\, t ,
  \qquad
  \frac{|\lambda-\lambda_n|}{\lambda_n}\;\lesssim\; C_\Omega\,\hat t .
\]
Here $C_\Omega$ is $O(1)$--$O(10)$ for non-pathological domains.
\textbf{Calibrate it once} per domain family against a known eigenvalue rather than
trusting a literature constant. Then set the design tolerance
\begin{equation}
  \boxed{\;
  \hat\epsilon=\frac{\epsilon}{10\,C_\Omega},
  \qquad
  \Lambda:=\ln(1/\hat\epsilon).
  \;}
  \label{eq:Lambda}
\end{equation}
The parameter $\Lambda$ is the single ``how many digits'' quantity that appears in
every count below. Design for the \emph{linear} Moler--Payne relation, not for the
quadratic behaviour that is sometimes observed empirically --- if the quadratic bonus
materialises, it is free margin.

\paragraph{Hard floor.} In IEEE double precision, $\sigma_{\min}$ of the GSVD pencil
plateaus around $10^{-13}$--$10^{-14}$ even with column equilibration, so
$\epsilon\lesssim 10^{-12}$ is the practical limit. Below that the remedy is extended
precision, not more basis functions.

\section{The organizing principle}
\label{sec:principle}

Two distinct things limit a particular-solution basis.

\begin{description}[leftmargin=2.2em,style=nextline]
\item[(a) Oscillation.]
  Any trace on $\partial\Omega$ has arclength wavenumber up to $\kappa$. The function
  $J_\nu(\kappa r)$ is exponentially small once $\nu>\kappa r$, with an Airy
  transition of width $\sim(\kappa r)^{1/3}$. So a corner expansion reaching radius
  $R$ needs orders up to $\nu\approx \kappa R + 2(\kappa R)^{1/3}$.
\item[(b) Analytic continuation.]
  A Fourier--Bessel series about $v_c$ converges like $(r/d_c)^{\nu}$, where $d_c$ is
  the distance from $v_c$ to the nearest \emph{singularity of the analytic
  continuation of $u$} --- \emph{not} the nearest corner. Beyond $r=d_c$ the series
  is useless no matter how many terms are added.
\end{description}

This is why single-corner MPS works spectacularly on the L-shape --- all the other
corners are $\pi/2$, hence regular, hence not singularities of the continuation ---
and fails on domains with several reentrant corners.

\begin{principle}[Partition of responsibility]
Each singular corner owns only the boundary arc within radius
\[
  R_c=\gamma\, d_c ,\qquad \gamma\in[0.3,0.5],
\]
and the MFS sources own everything else. Do not ask a corner expansion to cover the
whole boundary unless $d_c=\infty$.
\end{principle}

The \emph{obstruction set} used to compute $d_c$ is
\begin{align*}
  \Sstar=\ &\{\text{singular corners}\}\ \cup\
            \{\text{their reflections across each edge}\}\\
    &\cup\ \{\text{Schwarz-function singularities of curved arcs}\},
\end{align*}
with the practical proxy $2\times(\text{distance to the nearest non-adjacent edge})$
for the reflected images, and the radius of curvature for smooth curved arcs. Then
\[
  d_c=\dist\big(v_c,\ \Sstar\setminus\{v_c\}\big),
\]
capped by the largest radius actually needed,
$d_c\leftarrow\min\big(d_c,\ \max_{x\in\partial\Omega}|x-v_c|/\gamma\big)$.

There is a genuine optimum in $\gamma$: raising it shrinks the MFS workload but
inflates $M_c$ through both mechanisms. The total is quite flat over
$\gamma\in[0.3,0.5]$; use $\gamma=0.4$ unless you profile.

\section{Corner classification: which corners get enrichment?}
\label{sec:classification}

Compute $\mu_c=\dist(\alpha_c,\Z)$.

\begin{center}
\small
\begin{tabular}{@{}p{2.4cm}p{4.0cm}p{7.2cm}@{}}
\toprule
\textbf{Case} & \textbf{Condition} & \textbf{Treatment} \\
\midrule
Regular
  & $\mu_c=0$ (i.e.\ $\omega_c=\pi/m$)
  & No singular enrichment; $u$ extends by reflection. \\[3pt]
Weakly singular
  & $\mu_c>0$, $\alpha_c>\Lambda/(2\ln 15)$
  & Treat as regular. \\[3pt]
Singular
  & $\mu_c>0$, $\alpha_c\le\Lambda/(2\ln 15)$
  & Full corner-adapted treatment (\S\ref{sec:corner}). \\[3pt]
Near-integer
  & $0<\mu_c\lesssim0.05$
  & Singular, but scale the continuation term by
    $\min(1,\mu_c/0.05)$. \\
\bottomrule
\end{tabular}
\end{center}

The weakly singular cutoff comes from the observation that a smooth basis resolves
$r^{\alpha_c}$ to $\hat\epsilon$ with degree $\sim\hat\epsilon^{-1/(2\alpha_c)}$,
which is single digits when $\alpha_c$ is large. For $\hat\epsilon=10^{-11}$
($\Lambda\approx25$) the cutoff is $\alpha_c\lesssim4.6$: \textbf{corners wider than
about $39^\circ$ get corner-adapted terms; sharper spikes do not.}

\paragraph{Should regular corners receive corner-adapted terms?}
Not a $\Lambda$-sized budget. Two nuances:

\begin{itemize}[leftmargin=1.4em]
\item A \emph{small} set of $J_{jm}(\kappa r)\sin(jm\theta)$ at a regular corner
  ($m=\pi/\omega_c$ integral) is cheap and often improves conditioning, because it
  supplies well-localised angular content that the MFS represents only awkwardly. Use
  $M_c=\lceil \tfrac{\omega_c}{\pi}(\kappa R_c^{\mathrm{loc}}+5)\rceil$ with
  $R_c^{\mathrm{loc}}$ equal to half the shorter adjacent edge, and no continuation
  term. Treat this as optional.
\item The operationally important difference is in the MFS: \textbf{do not taper the
  source offset to zero at a regular corner} (\S\ref{sec:mfs}). That is where most of
  the savings live.
\end{itemize}

\section{Corner-adapted counts}
\label{sec:corner}

For each singular corner,
\begin{equation}
  \boxed{\;
  M_c=
  \underbrace{\Big\lceil \tfrac{\omega_c}{\pi}
    \big(\kappa_{\max}R_c+2(\kappa_{\max}R_c)^{1/3}+5\big)\Big\rceil}_{\text{oscillatory}}
  \;+\;
  \underbrace{\Big\lceil \tfrac{\omega_c}{\pi}\cdot
    \frac{\Lambda}{\ln(1/\gamma)}\Big\rceil}_{\text{continuation}}
  \;}
  \label{eq:Mc}
\end{equation}

Note the structure: the angle enters \emph{only} through the factor
$\omega_c/\pi=1/\alpha_c$, because the index spacing in $\nu$ is $\alpha_c$. Wider
corners cost proportionally more terms for the same $\nu_{\max}$ --- a $3\pi/2$
reentrant corner costs three times what a $\pi/2$ corner would at equal reach. That
single factor is the whole angle dependence; there is no separate ``reentrant corners
are harder'' term, because the difficulty is already captured by $1/\alpha_c$ and by
$d_c$ shrinking when singular corners cluster.

\paragraph{Conditioning cap.} Truncate at
\[
  \nu_{\max}^{(c)}\le \kappa R_c+\frac{\ln(1/\epsilon_{\mathrm{mach}})}{\ln(R_c/r_{\mathrm{mid}})},
  \qquad
  r_{\mathrm{mid}}=\text{median }|x-v_c|\text{ over the owned arc},
\]
since beyond that the column's dynamic range over the sampled boundary exceeds what
double precision can carry, and the GSVD will discard it anyway. Always
\textbf{column-equilibrate} (unit norm with respect to the combined
boundary\,+\,interior quadrature) before forming the pencil; without it the $J_\nu$
and $Y_k$ columns differ by hundreds of orders of magnitude.

\section{Fundamental-solution placement}
\label{sec:mfs}

Treat the offset sources and the lightning clusters as \emph{one graded curve}, not
two families. That is the cleanest way to keep them consistent.

\subsection{Local offset}

For $s\in\partial\Omega$,
\[
  \delta(s)=\min\Big(\delta_{\mathrm{amb}},\ \eta\cdot\dist(s,\Sstar)\Big),
  \qquad \eta\approx0.3,
  \qquad
  \delta_{\mathrm{amb}}=\min\big(0.25D,\ \tfrac{1}{3}\rho_{\mathrm{curv}}^{\min}\big).
\]
Sources go at $s+\delta(s)n(s)$ (outward normal). Because $\Sstar$ contains only
\emph{singular} corners and their images, $\delta$ stays at $\delta_{\mathrm{amb}}$
across regular corners and collapses only at singular ones --- exactly the geometric
clustering one wants, derived rather than imposed.

\subsection{Local spacing}

The classical MFS/Fourier result for a source curve at distance $\delta$ with spacing
$h$ is boundary error $\sim e^{-\pi\delta/h}$. Hence
\begin{equation}
  \boxed{\;
  h(s)=\min\Big(\frac{2\pi}{3\kappa_{\max}},\ \frac{\pi\,\delta(s)}{\Lambda}\Big)
  \;}
  \label{eq:h}
\end{equation}
The first argument is the Nyquist / three-points-per-wavelength floor; the second is
the accuracy requirement. For high-accuracy work the second almost always binds,
meaning \textbf{MFS sources are oversampled several-fold relative to Nyquist} --- that
is normal, and is the price of the offset.

\subsection{Total count}

\[
  N_{\mathrm{mfs}}\;\approx\;
  \max\Big(\frac{3\kappa_{\max}L}{2\pi},\ \frac{\Lambda L}{\pi\,\delta_{\mathrm{amb}}}\Big)
  \;+\;\sum_{c\ \mathrm{sing}} n_c^{\mathrm{clus}}.
\]
Since $h(s)\propto\dist(s,v_c)$ near a singular corner, the tapered region contributes
geometrically clustered sources with ratio $1+\pi\eta/\Lambda$ per step; the count for
the tail is
$n_c^{\mathrm{clus}}\approx \ln(\delta_{\mathrm{amb}}/s_{\min})
\big/\ln(1+\pi\eta/\Lambda)$.

\subsection{Where to stop the taper}

This is where the Fourier--Bessel basis and the clusters must not duplicate each other.

\begin{itemize}[leftmargin=1.4em]
\item \textbf{Corner has FB terms} (\S\ref{sec:corner}): stop the taper at
  $s_{\min}\approx 0.05\,R_c$ and keep only a short tail,
  $n_c^{\mathrm{clus}}\approx 8$--$12$ poles on the outward angle bisector. Their job
  is to bridge the region where the FB series is still converging but the ambient MFS
  cannot reach --- a conditioning bridge, not an approximation tool. More than this is
  wasted and actively hurts $\sigma_{\min}$.
\item \textbf{Corner has no FB terms} (weakly singular by \S\ref{sec:classification},
  or $M_c$ capped for conditioning): use full Gopal--Trefethen tapered clustering on
  the outward bisector,
  \[
    s_j=\delta_{\mathrm{amb}}\,e^{-\sigma(\sqrt{n_c}-\sqrt{j})},
    \qquad \sigma\approx4,\quad j=1,\dots,n_c,
  \]
  with root-exponential convergence $\sim e^{-\pi\sqrt{2\alpha_c n_c}}$, giving
  \begin{equation}
    \boxed{\;n_c^{\mathrm{clus}}\approx\frac{\Lambda^2}{2\pi^2\alpha_c}\;}
    \label{eq:lightning}
  \end{equation}
\end{itemize}

\subsection{Multipole orders}

Use plain monopoles ($Y_0$) for the clustered near-corner sources --- they are close to
the boundary and low order suffices. Reserve multipole stacks for a handful of far
centres, where they permit far fewer sources. For a centre at distance $\rho_j$ from
the arc it serves, of half-extent $a_j$,
\[
  K_j\approx \kappa\sqrt{\rho_j^2+a_j^2}
    +\frac{\Lambda}{\ln\!\big(\sqrt{\rho_j^2+a_j^2}/\rho_j\big)},
  \qquad
  K_j\le \frac{\ln(1/\epsilon_{\mathrm{mach}})}{\ln(1+D/\rho_j)} .
\]
The second inequality is a hard conditioning cap:
$Y_k(\kappa r)\sim -\frac{(k-1)!}{\pi}(2/\kappa r)^k$ has dynamic range
$(r_{\max}/r_{\min})^{K}$ across $\Omega$, and once that exceeds
$1/\epsilon_{\mathrm{mach}}$ the column carries no information.

Use only $Y_k$ at exterior centres, never $J_k$ --- the regular part about an exterior
centre is already spanned by the interior and corner Fourier--Bessel functions and
merely duplicates them.

\section{Collocation, the \texorpdfstring{$A_I$}{A\_I} block, and quadrature}
\label{sec:colloc}

\paragraph{Replace the interior points.} With interior sampling, $M_I$ must scale like
$A\kappa^2$ to control $\|u\|_{L^2(\Omega)}$ reliably, and undersampling produces
spurious tension minima that look exactly like eigenvalues. Use the Rellich identity
instead: for a Dirichlet solution of $\Delta u+\lambda u=0$,
\begin{equation}
  \|u\|_{L^2(\Omega)}^2
  =\frac{1}{2\lambda}\int_{\partial\Omega}(x\cdot n)\,|\partial_n u|^2\,ds ,
  \label{eq:rellich}
\end{equation}
with $x$ measured from any origin --- choose one making $x\cdot n>0$, i.e.\ a star
centre, if $\Omega$ is star-shaped. Then \emph{both} blocks of the pencil are boundary
matrices, $M_I$ becomes $O(\kappa)$ rather than $O(\kappa^2)$, and the
norm-equivalence constants relating $\sigma_{\min}$ to the true tension become
$\kappa$-independent --- which matters, since the whole $\epsilon\mapsto\hat\epsilon$
chain of \S\ref{sec:notation} assumes it.

\paragraph{Boundary nodes.} Take $M_B\approx 3N$ in total, with weights folded in as
$\sqrt{w_i}$ so the discrete norms actually approximate $L^2$; otherwise the
topological-equivalence constant drifts with $\kappa$ and the calibration of
$C_\Omega$ is meaningless. Use composite Gauss--Legendre panels, geometrically refined
toward each singular corner with ratio $1/2$, down to $\approx s_{\min}$ from
\S\ref{sec:mfs} --- roughly $\log_2(\ell_e/s_{\min})$ levels per adjacent edge.

\section{Total budget and precision floors}
\label{sec:budget}

\[
  N_{\mathrm{tot}}=\sum_{c\ \mathrm{sing}}M_c
    +\sum_{c\ \mathrm{reg}}M_c^{\mathrm{opt}}
    +N_{\mathrm{mfs}}.
\]

Sanity checks before building the matrix:

\begin{itemize}[leftmargin=1.4em]
\item If $N_{\mathrm{tot}}\gtrsim 600$ in double precision, conditioning will limit
  the GSVD before approximation power does. Relax $\epsilon$, reduce $\gamma$, or move
  to extended precision; adding functions will make $\sigma_{\min}$ \emph{worse}.
\item $N_{\mathrm{tot}}$ should scale like
  $O(\kappa_{\max}L)+O(\Lambda\cdot n_{\mathrm{sing}})$, i.e.\ like
  $\sqrt{\lambda_{\max}}$, not $\lambda_{\max}$. If the heuristic produces
  $O(\lambda)$ growth, something is wrong --- usually interior points.
\item If $\sigma_{\min}$ stops decreasing as $N$ increases, you are on the floor;
  enrichment will not help.
\end{itemize}

\section{\texorpdfstring{$\lambda$}{lambda}-sweep bookkeeping}
\label{sec:sweep}

Every count above depends on $\kappa$. \textbf{Do not} recompute the basis
continuously during a sweep: discontinuous basis changes put kinks in
$\sigma_{\min}(\lambda)$ and wreck root-finding and minimum tracking. Freeze the basis
on dyadic windows $\lambda\in[\lambda_0,2\lambda_0)$, sized at the window's upper end,
with hysteresis at the window boundaries. Within a window $\sigma_{\min}(\lambda)$ is
then smooth and the usual scaling and minimisation machinery applies.

\section{Adaptive correction loop}
\label{sec:adaptive}

The heuristic gets within a factor of roughly $1.5$; close the loop.

\begin{enumerate}[leftmargin=1.8em]
\item Build with the counts above; compute $\sigma_{\min}$ at a converged eigenvalue.
\item If $\hat\epsilon$ is not met, evaluate the boundary residual $|u(s)|$ panel by
  panel.
\item Attribute the largest residual to its owner: inside $R_c$ of a singular corner,
  increase $M_c$ by $25\%$; outside, halve $h$ on that arc, or --- if $\delta$ there is
  already at $\eta\dist(s,\Sstar)$ --- add a taper level.
\item If the residual is \emph{uniform} along $\partial\Omega$, you are
  conditioning-limited, not resolution-limited: increase $\delta_{\mathrm{amb}}$
  \emph{and} decrease $h$ together (keeping $\pi\delta/h=\Lambda$), or drop the weakest
  GSVD directions.
\end{enumerate}

\section{Worked example: the L-shape}
\label{sec:example}

Take $\lambda\approx100$, $\epsilon=10^{-10}$; $A=3$, $L=8$, $D=\sqrt5$, $\kappa=10$.
With $C_\Omega=10$ we get $\hat\epsilon=10^{-11}$ and $\Lambda=25.3$. The corners are
one at $3\pi/2$ ($\alpha=2/3$, singular) and five at $\pi/2$ (regular).

\begin{itemize}[leftmargin=1.4em]
\item Only one singular corner, so $\Sstar$ reduces to its own reflections: the
  nearest non-adjacent edge is at distance $1$, giving $d_c=2$. With $\gamma=0.4$,
  $R_c=0.8$.
\item Oscillatory term: $\kappa R_c=8$, so $\nu=8+4.0+5=17$ and
  $M^{\mathrm{osc}}=17\times1.5=26$.
\item Continuation term: $\nu=25.3/\ln 2.5=27.6$, so $M^{\mathrm{cont}}=41$. Hence
  $M_c=67$.
\item MFS: $\delta_{\mathrm{amb}}=0.25D=0.56$, so $h=\pi(0.56)/25.3=0.070$ --- the
  Nyquist floor would be $0.21$, so accuracy binds. Then
  $N_{\mathrm{mfs}}\approx 8/0.070\approx115$, plus a $\sim10$-pole bridge at the
  reentrant corner.
\item Regular corners: no taper and no continuation budget --- $\delta$ stays at
  $0.56$ all the way around them.
\end{itemize}

\noindent
Totals: $N_{\mathrm{tot}}\approx 195$, $M_B\approx600$ graded boundary nodes, no
interior points. That is the right order for ten digits at $\lambda\approx100$ on this
domain.

\section{Things to verify in a given implementation}
\label{sec:caveats}

\begin{description}[leftmargin=2.2em,style=nextline]
\item[Source-curve resonance.]
  Real MFS bases built from $Y_k$ on a closed exterior curve $\Gamma'$ can lose
  completeness at isolated $\lambda$ related to the spectrum of the region bounded by
  $\Gamma'$. It manifests as an isolated spike in $\sigma_{\min}$ that moves when
  $\Gamma'$ is perturbed. Cheap diagnostic: re-run suspicious $\lambda$ with
  $\delta_{\mathrm{amb}}$ scaled by $1.3$; genuine eigenvalues do not move. Including
  multipole orders mitigates this substantially.
\item[$C_\Omega$ calibration]
  is the weakest link in the chain --- but everything downstream is
  $\log(1/\hat\epsilon)$, so an order-of-magnitude error in $C_\Omega$ costs only
  about $10\%$ more basis functions. Do not over-engineer it.
\item[Degeneracies.]
  With multiplicity greater than one, a single $\sigma_{\min}$ is not enough: track
  the $m$ smallest generalized singular values and use the tension of the whole
  subspace. The counts above are unchanged.
\item[Tuned versus derived constants.]
  The values $\gamma=0.4$, $\eta=0.3$, $\sigma=4$, $M_B/N=3$, and the
  $2(\kappa R)^{1/3}+5$ transition margin are tuned, not derived. The structures
  $\Lambda/\ln(1/\gamma)$, $\pi\delta/h$, $\Lambda^2/(2\pi^2\alpha)$, and
  $\omega_c/\pi$ \emph{are} derived and should be trusted more than the constants in
  front of them.
\end{description}

\end{document}
